We describe the differential by composing a right flow module with the fixed flow manifolds. All flow manifolds, module manifolds, and bordisms below are compact; all boundary identifications include their higher intersections and compatible normal framings.

The geometric input is relative Pontryagin--Thom on finite diagrams of faces. We give the rounding and collapse comparison first, then realize algebraic extensions by relative transversality. Only framed bordism classes are needed; we do not assert a classification of flow bimodules or contractibility of the space of framings.

\begin{thm}\label{cor:SSofFFCat}
    Let $\fC$ be a flow category with finitely many objects in each degree and compact morphism manifolds, equipped with compatible neat embeddings and stable normal framings. Let $K_{\fC}$ be its associated coherent chain complex and let $(E_r^{n,p},d_r)$ be the spectral sequence of $\cI K_{\fC}$. Then
    \[
        E_1^{n,p}=\bigoplus_{\norm{x}=p}\Omega^{\fr}_{n-p}.
    \]
    Fix integers $n\geq p$ and $r\geq0$, and let
    $L=(L_x)_{\norm{x}=p}$ be closed framed $(n-p)$-manifolds representing $\ell\in E_1^{n,p}$.
    \begin{enumerate}
        \item The element $\ell$ survives to $E_{r+1}$, in the sense of \cref{def:survive}, if and only if $L$ extends to a degree-$n$ neatly embedded, normally framed right flow module $W$ over $\fC_{[p,p-r]}$.
        \item Given such $W$, for each object $z$ of degree $p-r-1$ form the closed framed $(n-p+r)$-manifold
        \[
            T_z(W)=M_{-,z}\circ W
            =\operatorname{round}\left(
              \bigcup_{\substack{x\in\fC\\p\geq\norm{x}\geq p-r}}
              W_x\times M_{x,z}\right).
        \]
        The union denotes gluing along the product faces of \cref{def:CompCat}. The collection
        \[
            \tau(W)=([T_z(W)])_{\norm{z}=p-r-1}
            \in E_1^{n-1,p-r-1}
        \]
        lies in $\cE_{r+1}^{n-1,p-r-1}$ and represents
        \[
            d_{r+1}^{n,p}([\ell]_{r+1})
            =[\tau(W)]_{r+1}.
        \]
        This page class is independent of $W$; the class $\tau(W)$ in $E_1$ need not be.
        \item The differential $d_{r+1}^{n,p}([\ell]_{r+1})$ vanishes if and only if there is a choice of $W$ extending $L$ for which every $T_z(W)$ is framed null-bordant. For that choice, null-bordisms with the prescribed corner faces extend $W$ over $\fC_{[p,p-r-1]}$.
    \end{enumerate}
\end{thm}

The constructions use only bounded grading windows, which are finite under the stated hypothesis. No new unbounded realization construction is asserted here: $K_{\fC}$ is the given coherent chain complex, and the argument applies to its finite restrictions.

\subsection{Gluing a right and a left module}

\begin{defn}\label{def:CompCat}
    Let $a\leq b\leq n$, let $\fC$ be a finite framed flow category supported in $[b,a]$, and let $W,V$ be compact framed right and left flow modules over $\fC$ of degrees $n$ and $a$, respectively. Set $I=\{a,\ldots,b\}$. For a nonempty subset $S=\{i_0>\cdots>i_k\}\subseteq I$, define
    \[
        D_S(W,V)=
        \bigsqcup_{\norm{x_j}=i_j}
        W_{x_0}\times M_{x_0,x_1}\times\cdots
        \times M_{x_{k-1},x_k}\times V_{x_k}.
    \]
    For $k=0$ this means $\bigsqcup_{\norm{x}=i_0}W_x\times V_x$. Flow composition and the module actions give a diagram indexed by nonempty subsets with arrows $T\to S$ for $S\subseteq T$. Write
    \[
        |D(W,V)|=\colim_{\varnothing\ne S\subseteq I}D_S(W,V).
    \]
    We denote its framed rounding, constructed in \cref{lem:colimIsMfld}, by $V\circ W$.
\end{defn}

\begin{lem}[The manifold and its framing]\label{lem:colimIsMfld}
    For the modules of \cref{def:CompCat}, $|D(W,V)|$ admits a rounding to a closed smooth manifold of dimension $n-a$. The compatible normal framings determine a stable framing of this rounding. Its framed bordism class is independent of the compatible collars and rounding choices.
\end{lem}
\begin{proof}[Proof sketch]
    Put $d=n-a$. The stratum $D_S$ has dimension $d+1-|S|$. Boundary exhaustion and associativity give, near its interior, the model
    \[
        \bbR^{d+1-|S|}\times\partial\Hlf^{|S|}.
    \]
    Indeed, the $|S|$ top pieces are its incident coordinate faces, and their intersections are exactly the products indexed by supersets of $S$. Thus the colimit has no unpaired boundary.

    Choose product collars, compatibly from the deepest strata outwards. Relative collars and framed composition are supplied by \cite[Lemmas~3.19, 4.9, and 4.41--4.43]{porcelli2024bordism} (arXiv version~3). In that construction the pieces $W_x\times V_x$ acquire an extra collar direction. On a double overlap the two collar coordinates are exchanged; their zero faces are precisely the identifications in $D(W,V)$. Associativity gives the same description on higher overlaps. A nearby regular hypersurface therefore rounds this particular colimit.

    For our normal-framing convention, stabilize the compatible product embeddings into the coordinate faces of a common orthant. On an overlap both normal trivializations restrict to the same ordered product trivialization on $D_S$. The exchanged collar coordinates are tangent to the thickened ambient orthant, not additional normal-framing factors. Rounding these coordinates in both the ambient space and the submanifold preserves that normal trivialization. Restricting to the rounded ambient boundary gives the asserted framed manifold. Applying the construction to a family gives a framed bordism between choices. The framing data are held fixed throughout.
\end{proof}

\begin{prop}[Collapse and composition]\label{prop:flow_PT_composition}
    Let $\fC,W,V$ be as in \cref{def:CompCat}, with compatible neat embeddings and normal framings. Their collapse data determine morphisms
    \[
        f_W:(\Sigma^{n-b}\bbS)_b\longrightarrow K_{\fC},
        \qquad
        g_V:K_{\fC}\longrightarrow\bbS_a
    \]
    in $\Ch_{[b,a]}(\Sp)$. Under \cref{prop:ChModels.composition}, their composite corresponds to
    \[
        \operatorname{PT}(V\circ W):
        \Sigma^{n-a}\bbS\longrightarrow\bbS.
    \]
\end{prop}
\begin{proof}[Proof sketch]
    Choose product collars and tubular neighborhoods by induction over faces, retaining the choices already made on their intersections. Extend a product metric from a smaller boundary collar and use a normal exponential neighborhood; compactness permits a common smaller radius. Relative neat embeddings exist after stabilization \cite[Proposition~2.8]{Genauer12}; compare \cite[Appendix~A]{cote2023equivariant}. There are finitely many pieces, so these choices give compatible finite-dimensional collapse maps. The module boundary identities are the enriched endpoint identities of \cref{prop:ChModels.endpoint_maps}. Keep these maps in their suspended form, apply a free orthogonal spectrum functor, and use \cref{lem:ChModels.transport} to pass to the fibrant representatives used for derived mapping spaces.

    Here is the local comparison before replacement. For nonempty $S\subseteq I$, embed $D_S$ in
    \[
        E_S=\bbR^N\times\Hlf^{I\setminus S}.
    \]
    Its normal bundle $\nu_S$ has rank $c=N+|I|-(n-a)-1$, independent of $S$. Choose tubular neighborhoods which are products on all collars. Restriction to a face then identifies both the normal bundles and their collapse maps, giving a diagram
    \[
        E_S^+\longrightarrow\Th(\nu_S).
    \]
    The inclusions of unions of boundary faces are cofibrations, so the finite face colimits compute homotopy colimits. The ambient colimit is
    $\bbR^N\times\partial\Hlf^I$, whose compactification is $S^{N+|I|-1}$.
    On each collar the embedding and tubular neighborhood are products, so the transition maps of the normal bundles restrict to the same maps on every intersection. They assemble to a bundle $\nu$ over $|D(W,V)|$, and
    \[
        \colim_S\Sigma^\infty\Th(\nu_S)
        \simeq\Th_{\Sp}(\nu).
    \]
    This target identification can also be obtained from preservation of colimits by the Thom-spectrum functor \cite[Section~3.3]{ABGHR14Thom}, after regarding the bundle data as a diagram of stable spherical fibrations.

    In a product tubular chart, rounding changes only the base collar coordinates and leaves the normal coordinate unchanged. Outside these charts every collapse sends the complement of the tubular neighborhood to the basepoint. Consequently these local identifications give a homotopy from the assembled collapse to the collapse of the rounded embedding, relative to the unchanged regions. After using the framing to map the Thom space to $S^c$, its restriction to the $i$-face is the product collapse for $\bigsqcup_{\norm{x}=i}W_x\times V_x$. These are exactly the face maps in \cref{eq:ChModels.face_composite}; the coend relations identify their common restrictions. Thus the assembled map is the categorical composite.
\end{proof}
\verifyproof{Convention check remaining from G2: the normal-framed gluing and coend maps agree with the same ordered boundary-sphere identification. Check its sign against the independently chosen suspension/connecting-morphism identifications in \cref{lem:Ch_diff}. The one-degree case must give $[L\times M]$ with the product framing. This is a sign comparison, not a remaining existence problem for collars or collapse maps.}

\subsection{Realizing extensions and restoring corner faces}

\begin{lem}[Realization with a fixed top representative]\label{lem:flow_relative_realization}
    Let $\fC$ be a finite neatly embedded, normally framed flow category supported in $[p,a]$, let $n\geq p$, and let
    \[
        f:(\Sigma^{n-p}\bbS)_p\longrightarrow K_{\fC}
    \]
    be a morphism of coherent chain complexes. Suppose its component in degree $p$ is represented by the Pontryagin--Thom maps of closed framed manifolds $(L_x)_{\norm{x}=p}$. After stabilization, $f$ is represented by a compact neatly embedded, normally framed right flow module $W$ of degree $n$ with $W_x=L_x$ in degree $p$.
\end{lem}
\begin{proof}[Proof sketch]
    Use the finite based-space collapse diagram of \cref{con:topChOfCat}, retaining its fixed flow structure maps. By \cref{lem:ChModels.finite_representatives}, $f$ has simultaneous unstable endpoint representatives after suspension. The relative assertion of that lemma lets us keep the prescribed top collapse maps of $L$, rather than merely their stable classes.

    Proceed downwards in degree. Suppose the module pieces above $i$ have been constructed. Their products with $M_{y,x}$, for $\norm{x}=i$, prescribe the collapse on all boundary faces of the next orthant. These prescriptions agree on intersections by the module identities already constructed. The attaching pair is
    \[
       (\partial\Hlf^{p-i})_+\wedge S^{N+n-p}
       \ \subseteq\
       (\Hlf^{p-i})_+\wedge S^{N+n-p}.
    \]
    Its inclusion is a cofibration. Replace the boundary restriction of the chosen representative by the prescribed product collapse using a relative homotopy, and extend that homotopy over the attaching cell. Make the map a product near the boundary. Smooth it near the preimages of chosen non-basepoints in the finitely many target spheres, and perturb it to be transverse there, keeping the boundary collars fixed. The corner transversality and inverse-image construction are those of \cite[Section~9 and Theorem~3.17]{Genauer12}, applied relative to the already transverse collar.

    The inverse image in the $x$-sphere defines $W_x$. Its dimension is
    $(N+n-p)+(p-i)-N=n-i$; its boundary faces are exactly $W_y\times M_{y,x}$. The differential of the map identifies its normal bundle with the fixed tangent space of the target sphere, giving the required normal framing. The relative Pontryagin--Thom homotopy replaces the map by the collapse of this inverse image without changing it on the boundary collars. Extend this homotopy over the remaining cells of $Q_X$ before proceeding to degree $i-1$. Thus each step preserves the entire endpoint morphism $f$, not just its top component. Finally choose the embeddings in the ordered product coordinates of \cref{def:FramedFlowModule}, inductively relative to the prescribed boundary collars. Stabilized relative embedding and its parametric version identify these with the embeddings just constructed, carrying their normal framings. Unused face coordinates above $p$ can be held strictly positive. The finite induction constructs $W$ with $W_x=L_x$ in top degree.
\end{proof}

\begin{lem}[Restoring the faces of a null-bordism]\label{lem:flow_restore_corners}
    Let $W$ be a degree-$n$ framed right flow module over $\fC_{[p,a]}$ and let $z$ have degree $a-1$. Suppose the framed rounding $T_z(W)$ is null-bordant. Then a framed null-bordism can be given corners and, after stabilization, a neat embedding extending the prescribed boundary data, so that it defines a manifold $W_z$ of dimension $n-a+1$ with faces
    \[
        \partial_i W_z
        =\bigsqcup_{\norm{x}=i}W_x\times M_{x,z},
        \qquad a\leq i\leq p,
    \]
    and with the higher intersections prescribed by flow composition and the module action. Performing this construction for all $z$ of degree $a-1$ extends $W$ over $\fC_{[p,a-1]}$.
\end{lem}
\begin{proof}[Proof sketch]
    Let $B_z$ be a framed null-bordism together with its framed boundary identification $\partial B_z\cong T_z(W)$. Retain the collar region between the original cornered union and its rounding in \cref{lem:colimIsMfld}. Attach its rounded end to a boundary collar of $B_z$. The other end has precisely the prescribed faces and their original intersections; no choices of new face identifications are involved.

    Use the outward-normal-first boundary convention. The two collars being attached have opposite outward normals, so the attaching map reverses that line and identifies the remaining framed directions by the given framed boundary identification. Extend over a product collar. This transports the stable framing of $B_z$ to the manifold with restored corners and agrees with every prescribed product framing. Choose a neat embedding relative to this entire boundary collar after stabilization \cite[Proposition~2.8]{Genauer12}. Its stable normal framing is the complement of the transported stable tangent framing; a further stabilization realizes the prescribed boundary homotopy of normal trivializations on a collar. Distinct objects in degree $a-1$ impose no additional identities on one another.
\end{proof}

\subsection{Proof of the differential theorem}

\begin{proof}[Proof sketch of \cref{cor:SSofFFCat}]
    The degree-$p$ spectrum is a wedge of sphere spectra, giving the stated first page by Pontryagin--Thom.

    If $L$ extends to $W$ over $\fC_{[p,p-r]}$, its collapse data give
    \[
        f_W:(\Sigma^{n-p}\bbS)_p
        \longrightarrow K_{\fC_{[p,p-r]}}
    \]
    extending the class of $L$. Proposition~\ref{lem:ChSur} implies survival to $E_{r+1}$. Conversely, that proposition supplies an algebraic extension for every surviving class; \cref{lem:flow_relative_realization} realizes it with the prescribed top manifolds $L$.

    Put $a=p-r$. For each $z$ of degree $a-1$, the manifolds $M_{x,z}$ form a left module of degree $a$ over $\fC_{[p,a]}$. Its collapse data are the lower endpoint data of $K_{\fC}$ at $z$. By \cref{lem:Ch_diff}, composing with $f_W$ gives the $z$-component of the raw obstruction $\partial_{r+1}(f_W)$. Proposition~\ref{prop:flow_PT_composition} identifies this component with $[T_z(W)]$. The same algebraic differential description shows that this collection belongs to $\cE_{r+1}$ and represents $d_{r+1}([\ell]_{r+1})$.

    Finally suppose this page differential is zero. Lemma~\ref{lem:zero_dr_extension}, applied with $r+1$, supplies an extension $f'$ of the same top class for which
    \[
        \partial_{r+1}(f')=0
        \quad\text{in }E_1^{n-1,p-r-1}.
    \]
    Realize $f'$ by a module $W'$ using \cref{lem:flow_relative_realization}. Its Toda manifolds are framed null-bordant by the composition comparison. Apply \cref{lem:flow_restore_corners} to extend $W'$ one degree further. Conversely, null-bordant Toda manifolds give zero raw obstruction and hence zero page differential. The replacement of $W$ by $W'$ is essential: vanishing on the page does not assert null-bordism for every initial choice of $W$.
\end{proof}

\addtopaper{Add a direct geometric description of \(E_r^{n,p}\): it is the group of level-\(p\) framed collections \(W_p\) that extend to degree-\(n\) right flow modules over \(\fC_{[p,p-r+1]}\), modulo the level-\(p\) classes obtained by gluing degree-\((n+1)\) right flow modules \(V\) over the upper window \(\fC_{[p+r-1,p+1]}\) to the fixed flow manifolds. At an object \(x\) of index \(p\), this glued collection is \(\bigcup_{p+1\leq |y|\leq p+r-1}V_y\times M(y,x)\); these classes encode the cumulative incoming-boundary relations through \(d_{r-1}\).}
